\section{Overview}

This chapter reviews four strands of literature relevant to this dissertation: the growing use of large language models (LLMs) in medical and mental health contexts; documented bias and fairness concerns in AI applied to health; geographic, cultural, and cross-lingual variation in LLM behaviour specifically; and existing methodological approaches to auditing LLMs for such variation, including the limitations of automated content coding. The chapter closes by identifying the gap this dissertation addresses.

\section{LLMs in Medical and Mental Health Contexts}

Large language models have demonstrated substantial clinical knowledge and, in several evaluations, response quality that patients and even physicians can find comparable to or preferable over human-authored answers. \citet{singhal2023clinical} show that general-purpose LLMs can encode and apply clinical knowledge across a range of medical question-answering benchmarks, while \citet{ayers2023chatbot} find that, in a direct comparison on a public social media forum, a chatbot's responses to patient questions were rated by evaluators as more empathetic than physicians' responses in the majority of cases. \citet{wang2024healthcare} and \citet{busch2025patientcare} both survey the rapidly expanding set of applications of LLMs across medical and patient-care settings, while noting that evaluation methodology has lagged behind deployment.

Within mental health specifically, \citet{hua2025mentalhealth} provide a scoping review of LLM applications in mental health care, and \citet{mesquiti2026mentalhealth} review the use of LLMs in mental health research and treatment more broadly, both concluding that the evidence base for safety and effectiveness remains thin relative to the pace of adoption. \citet{rousmaniere2025mentalhealth} document patterns of real-world use of LLMs as mental health resources in the United States, providing direct evidence that people already use general-purpose LLMs -- not purpose-built clinical tools -- for this purpose. \citet{saha2025linguistic} conduct a linguistic comparison of AI- and human-written responses to online mental health queries, finding systematic stylistic and structural differences that may affect how such responses are received.

Two contributions are particularly relevant to the evaluation methodology used in this dissertation. \citet{golden2024faita} propose the Framework for AI Tool Assessment in Mental Health (FAITA-Mental Health), a structured scale for evaluating AI-powered mental health tools, reflecting a broader move toward more rigorous, multi-dimensional assessment frameworks rather than single-axis accuracy metrics. \citet{santos2026clinical} evaluate the clinical safety of LLMs specifically in response to high-risk mental health disclosures, a safety-critical scenario directly analogous to several of the scenarios used in this dissertation's prompt battery (Appendix A).

\section{Bias, Fairness, and Equity in AI and Healthcare}

A substantial literature documents that AI systems applied to healthcare can encode and reproduce demographic disparities. \citet{omar2024disparities} conduct a systematic review of demographic disparities in medical LLMs, finding evidence of differential performance and recommendation quality across race, gender, and socioeconomic groups. \citet{li2024inequities} similarly review health inequities arising from AI more broadly, situating algorithmic bias within the wider literature on structural inequity in healthcare access and delivery. \citet{bhanot2021fairness} examine the specific problem of fairness in synthetic healthcare data, a methodologically adjacent concern to the present dissertation's own construction of a synthetic (prompt-based) rather than naturalistic dataset.

Economic and health-services literature provides a complementary framing for the horizontal-equity concept this dissertation applies to geography rather than demography. \citet{zhao2020equity} propose a framework for measuring and decomposing horizontal inequity in the delivery of health care, and \citet{basar2020equity} apply a related equity lens to mental healthcare utilisation specifically, from a gender perspective, in the context of Turkey. Both works motivate the framing adopted here: that a fair system should provide comparable actionable support to users in comparable need, regardless of a characteristic -- in this dissertation's case, location -- that is not itself relevant to the clinical merit of the request. \citet{osborn2024students} provide a qualitative account, through the lens of candidacy theory, of how structural and perceptual barriers shape university students' access to mental health services in England, illustrating that geographic and institutional context shapes not only whether support exists but whether it is perceived as reachable -- a distinction this dissertation's localisation measure (Chapter 3) attempts to operationalise for LLM responses.

\section{Geographic, Cultural, and Cross-Lingual Variation in LLMs}

Closest to the present dissertation's central question is a body of work documenting that LLM outputs vary systematically with geographic or cultural framing, independent of the substantive content of the query. \citet{decoupes2025geographical} evaluate geographical distortions in language models directly, providing methodological precedent for treating location as an experimental variable in LLM evaluation. \citet{lyu2024regional} document regional bias specifically in monolingual English language models, showing that models trained predominantly on data from certain English-speaking regions perform unevenly when queried about or from other regions. \citet{jin2024english} show that LLMs answering healthcare queries perform measurably better when queried in English than in other languages, even holding the substantive query constant -- a finding with direct relevance to users in non-Anglophone, lower-income settings who may be more likely to query in a local language, and with relevance to this dissertation's choice to hold language constant while varying only location. \citet{orosoo2024multilingual} examine multilingual chatbot performance in cross-cultural communication more broadly.

Two further contributions bear on the interpretive framing of cultural content in LLM output specifically. \citet{healey2024bias} propose methods for evaluating nuanced bias in free-response LLM answers, moving beyond simple keyword-based bias metrics toward more holistic judgements of the kind this dissertation's human validation protocol also attempts. \citet{gupta2024transparency} likewise focus on enhancing transparency and mitigating bias in LLM responses, with an emphasis on making bias patterns auditable rather than merely documenting their existence. Finally, \citet{hall2024representation} provides the foundational cultural-studies framing of representation and signifying practices used in this dissertation's interpretation of religious and cultural framing (H3): that a model's default invocation of a particular cultural or religious support structure for a given region, absent any user-stated preference, constitutes a representational choice with the same stereotyping risks documented in Hall's broader account of media representation.

\section{Auditing LLMs: Methodology and the Limits of Automated Coding}

Much of the empirical literature reviewed above, and much of the wider LLM-audit literature, relies on some form of automated content coding -- keyword dictionaries, embedding-based similarity, or LLM-assisted labelling -- to convert free-text model output into a quantitative outcome measure at scale. \citet{dai2023llm} demonstrate one such approach directly, using an LLM in the loop to conduct thematic analysis of qualitative data, while \citet{depaoli2024thematic} conduct a critical exploration of the limits of using LLMs to perform inductive thematic analysis of interview data, concluding that such automated coding requires careful validation against human judgement rather than being treated as a drop-in replacement for it. This concern is central to the present dissertation's methodology (Chapter 3): as demonstrated directly in Chapter 4, an automated outcome measure can pass a face-validity check, agree well with an alternative automated coding rule, and still measure something other than what it claims to measure, in a way detectable only by comparison against independent human labels.

Where embedding- or topic-based automated methods were considered as a complement to keyword-based coding during the design of this dissertation's pipeline, the underlying techniques -- sentence embeddings \citep{reimers2019sentencebert}, dimensionality reduction \citep{mcinnes2018umap}, density-based clustering \citep{campello2013density}, and class-based topic modelling \citep{grootendorst2022bertopic} -- are well established in the broader NLP literature, though their application to bias auditing specifically remains less mature than the keyword-based approaches critiqued above.

\section{Summary and Research Gap}

The literature establishes three things clearly: LLMs are used as de facto mental-health resources by real users \citep{rousmaniere2025mentalhealth, saha2025linguistic}; their outputs vary with geographic, linguistic and cultural framing in domains beyond mental health \citep{decoupes2025geographical, lyu2024regional, jin2024english}; and automated content-coding pipelines require comparison with human judgement \citep{dai2023llm, depaoli2024thematic}. This dissertation extends that work to mental-health guidance across a deliberately diverse city sample and reports its own validation failure transparently. The frozen revised actionability and localisation measures were evaluated on a fresh 160-response hold-out and did not meet the pre-specified acceptance criteria.
